% ==============================================================================
% CHƯƠNG 7: TASK 5 — ĐI DÂY PCB CÓ PHẠT BẺ GÓC
% ==============================================================================
\chapter{Task 5: Đi dây mạch in PCB có phạt bẻ góc}
\label{chap:task5}

\section{Mô hình hóa}
\label{sec:t5-mo-hinh}

Bảng mạch là lưới $m \times n$ (\ma{pcbGrid}), ô bằng $1$ là vùng bị chiếm (linh kiện), các ô
khác đi dây được. Dây dẫn chỉ đi theo bốn hướng trục (Up, Down, Left, Right) như trong định tuyến
mê cung kinh điển \cite{lee1961}. Mỗi bước dài một đơn vị; nếu hướng của bước khác hướng của bước
ngay trước thì cộng thêm phí bẻ góc $B = 0{,}5$:
\begin{equation}
  g(n_{k}) = g(n_{k-1}) + 1 + B \cdot [\,\text{dir}_k \ne \text{dir}_{k-1}\,], \qquad B = 0{,}5,
  \label{eq:t5-cost}
\end{equation}
với quy ước bước đầu tiên không có hướng trước nên không bị phạt. Heuristic: Manhattan
(\ma{mode} $1$) hoặc Chebyshev (\ma{mode} $2$) tới ô ra.

\section{Thuật toán biến đổi lưới thành đồ thị}
\label{sec:t5-bien-doi}

\paragraph{Đánh số đỉnh} Ô $(x, y)$ ứng với đỉnh $\text{id}(x, y) = x \cdot n + y$; ngược lại
$x = \lfloor \text{id} / n \rfloor$, $y = \text{id} \bmod n$. Ánh xạ này là song ánh từ tập ô lên
$\{0, \dots, mn - 1\}$.

\paragraph{Dựng cạnh} Với mỗi ô tự do $(x, y)$ và mỗi hướng trục $(\Delta x, \Delta y)$, nếu ô
$(x + \Delta x, y + \Delta y)$ nằm trong lưới và tự do thì thêm cung có trọng số $1$. Mỗi cạnh
vô hướng của lưới được sinh hai lần (một lần từ mỗi đầu), nên ma trận kết quả đối xứng.

\begin{lstlisting}[style=gia, caption={Biến đổi lưới PCB thành ma trận trọng số}, label={lst:t5-bien-doi}]
function BuildGraph(pcbGrid, $m$, $n$, weightMatrix)
    set every weightMatrix$[i][j] \gets 0$
    for $x \gets 0$ to $m-1$, $y \gets 0$ to $n-1$
        if pcbGrid$[x][y] = 1$ then continue
        $u \gets x \cdot n + y$
        for each $(\Delta x, \Delta y, d)$ in {Up, Down, Left, Right}
            $(x', y') \gets (x + \Delta x, y + \Delta y)$
            if $(x', y')$ inside grid and pcbGrid$[x'][y'] \ne 1$ then
                $v \gets x' \cdot n + y'$
                wires$[u]$.add($v$, $d$)
                if $u < 100$ and $v < 100$ then weightMatrix$[u][v] \gets 1$
\end{lstlisting}

Thuật toán duyệt mỗi ô và bốn hướng một lần nên tốn $O(mn)$ thời gian, cộng $O(100^2)$ để xóa
ma trận đầu vào. Ma trận \ma{weightMatrix[100][100]} chỉ chứa được lưới có $mn \le 100$ ô; với lưới
lớn hơn, danh sách kề nội bộ \ma{wires} vẫn đầy đủ nên việc tìm đường không bị ảnh hưởng.

\begin{figure}[H]
  \centering
  \begin{tikzpicture}[x=1.15cm, y=1.15cm]
    \foreach \x/\y/\t in {0/0/0,0/1/0,0/2/0,0/3/0, 1/0/0,1/1/1,1/2/1,1/3/0,
                          2/0/0,2/1/0,2/2/0,2/3/0, 3/0/1,3/1/0,3/2/0,3/3/0} {
      \pgfmathtruncatemacro{\id}{\x*4+\y}
      \ifnum\t=1 \fill[bkxanhdam!80] (\y, -\x) rectangle ++(1, -1);
      \else \fill[white] (\y, -\x) rectangle ++(1, -1); \fi
      \draw[bkxam!60] (\y, -\x) rectangle ++(1, -1);
      \node[font=\scriptsize, text=bkxam] at (\y+0.5, -\x-0.5) {\id};
    }
    \draw[do, line width=2pt, -{Stealth[length=3mm]}]
      (0.5, -0.5) -- (3.5, -0.5) -- (3.5, -3.5);
    \node[font=\small\bfseries, anchor=south] at (0.5, -0.45) {pin};
    \node[font=\small\bfseries, anchor=north] at (3.5, -3.55) {pad};
  \end{tikzpicture}
  \caption{Bảng mạch mẫu $4 \times 4$ (ô xanh đậm bị chiếm, số là chỉ số đỉnh) và đường dây tìm được}
  \label{fig:t5-luoi}
\end{figure}

\begin{table}[H]
  \centering
  \caption{Ma trận trọng số \ma{weightMatrix} được điền cho bảng mạch ở \Hinh{fig:t5-luoi}}
  \label{tab:t5-ma-tran}
  \footnotesize
  \setlength{\tabcolsep}{3.2pt}
  \begin{tabular}{c|cccccccccccccccc}
    & \textbf{0} & \textbf{1} & \textbf{2} & \textbf{3} & \textbf{4} & \textbf{5} & \textbf{6} & \textbf{7}
    & \textbf{8} & \textbf{9} & \textbf{10} & \textbf{11} & \textbf{12} & \textbf{13} & \textbf{14} & \textbf{15} \\
    \hline
    \textbf{0} & \mz & \mo & \mz & \mz & \mo & \mz & \mz & \mz & \mz & \mz & \mz & \mz & \mz & \mz & \mz & \mz \\
    \textbf{1} & \mo & \mz & \mo & \mz & \mz & \mz & \mz & \mz & \mz & \mz & \mz & \mz & \mz & \mz & \mz & \mz \\
    \textbf{2} & \mz & \mo & \mz & \mo & \mz & \mz & \mz & \mz & \mz & \mz & \mz & \mz & \mz & \mz & \mz & \mz \\
    \textbf{3} & \mz & \mz & \mo & \mz & \mz & \mz & \mz & \mo & \mz & \mz & \mz & \mz & \mz & \mz & \mz & \mz \\
    \textbf{4} & \mo & \mz & \mz & \mz & \mz & \mz & \mz & \mz & \mo & \mz & \mz & \mz & \mz & \mz & \mz & \mz \\
    \textbf{5} & \mz & \mz & \mz & \mz & \mz & \mz & \mz & \mz & \mz & \mz & \mz & \mz & \mz & \mz & \mz & \mz \\
    \textbf{6} & \mz & \mz & \mz & \mz & \mz & \mz & \mz & \mz & \mz & \mz & \mz & \mz & \mz & \mz & \mz & \mz \\
    \textbf{7} & \mz & \mz & \mz & \mo & \mz & \mz & \mz & \mz & \mz & \mz & \mz & \mo & \mz & \mz & \mz & \mz \\
    \textbf{8} & \mz & \mz & \mz & \mz & \mo & \mz & \mz & \mz & \mz & \mo & \mz & \mz & \mz & \mz & \mz & \mz \\
    \textbf{9} & \mz & \mz & \mz & \mz & \mz & \mz & \mz & \mz & \mo & \mz & \mo & \mz & \mz & \mo & \mz & \mz \\
    \textbf{10} & \mz & \mz & \mz & \mz & \mz & \mz & \mz & \mz & \mz & \mo & \mz & \mo & \mz & \mz & \mo & \mz \\
    \textbf{11} & \mz & \mz & \mz & \mz & \mz & \mz & \mz & \mo & \mz & \mz & \mo & \mz & \mz & \mz & \mz & \mo \\
    \textbf{12} & \mz & \mz & \mz & \mz & \mz & \mz & \mz & \mz & \mz & \mz & \mz & \mz & \mz & \mz & \mz & \mz \\
    \textbf{13} & \mz & \mz & \mz & \mz & \mz & \mz & \mz & \mz & \mz & \mo & \mz & \mz & \mz & \mz & \mo & \mz \\
    \textbf{14} & \mz & \mz & \mz & \mz & \mz & \mz & \mz & \mz & \mz & \mz & \mo & \mz & \mz & \mo & \mz & \mo \\
    \textbf{15} & \mz & \mz & \mz & \mz & \mz & \mz & \mz & \mz & \mz & \mz & \mz & \mo & \mz & \mz & \mo & \mz \\  \end{tabular}
\end{table}

Các hàng $5$, $6$, $12$ toàn $0$ vì đó là ô bị chiếm. Hàng $3$ (ô góc $(0,3)$) chỉ có hai số $1$ tại
đỉnh $2$ và $7$; hàng $10$ (ô $(2,2)$) có ba số $1$ tại $9, 11, 14$ vì láng giềng phía trên $(1,2)$ bị
chiếm. Ma trận đối xứng, đúng với tính vô hướng của lưới.

\section{Duy trì hướng trước đó trong A*}
\label{sec:t5-huong}

Phí bẻ góc phụ thuộc vào \emph{bước trước}, nên nếu chỉ dùng ô làm trạng thái thì chi phí của một
cung không còn cố định và A* (vốn giả định chi phí cung chỉ phụ thuộc hai đầu cung) có thể sai:
hai đường cùng tới một ô với hướng khác nhau có tương lai khác nhau, nhưng đường có $g$ lớn hơn
một chút có thể bị loại dù sau đó đi tiếp rẻ hơn.

Chương trình giải quyết bằng cách \nhanmanh{mở rộng không gian trạng thái}: mỗi trạng thái là
cặp $(\text{ô},\ \text{hướng vào})$ với hướng vào thuộc $\{\text{Up}, \text{Down}, \text{Left},
\text{Right}, \varnothing\}$, trong đó $\varnothing$ chỉ dùng cho điểm xuất phát. Trạng thái được
mã hóa thành số nguyên $5 \cdot \text{id}(\text{ô}) + \text{hướng}$, nên lõi A* dùng chung không cần
sửa gì.

\begin{khungcongthuc}[Cung trong đồ thị trạng thái]
  \begin{equation}
    (u, d_{\text{vào}}) \xrightarrow{\ \ d\ \ } (v, d), \qquad
    c = 1 + \begin{cases} 0{,}5 & \text{nếu } d_{\text{vào}} \ne \varnothing \text{ và } d_{\text{vào}} \ne d \\ 0 & \text{ngược lại} \end{cases}
    \label{eq:t5-state}
  \end{equation}
\end{khungcongthuc}

\begin{khungma}{Algo.cpp (trích findPCBPath)}{C++}
for (int incoming = 0; incoming < STATES_PER_CELL; ++incoming) {
    int state = u * STATES_PER_CELL + incoming;
    graph.h[state] = h;
    isGoal[state] = u == grid.id(exitX, exitY);
    for (size_t i = 0; i < wires[u].size(); ++i) {
        const Wire& wire = wires[u][i];
        bool bends = incoming != NO_DIRECTION && incoming != wire.direction;
        double cost = TRACE_UNIT + (bends ? BEND_PENALTY : 0.0);
        graph.addEdge(state, wire.to * STATES_PER_CELL + wire.direction, cost);
    }
}
\end{khungma}

\begin{dinhly}
  Đường đi trong đồ thị trạng thái từ $(s, \varnothing)$ tới một trạng thái bất kỳ của ô ra tương
  ứng một-một với đường dây trên lưới, và chi phí của chúng bằng nhau. Do đó đường ngắn nhất trong
  đồ thị trạng thái là đường dây tối ưu theo \eqref{eq:t5-cost}.
\end{dinhly}

\begin{loigiai}
  Một đường dây $s = c_0, c_1, \dots, c_k$ với các hướng $d_1, \dots, d_k$ ứng với dãy trạng thái
  $(c_0, \varnothing), (c_1, d_1), \dots, (c_k, d_k)$; ngược lại bỏ thành phần hướng của dãy trạng
  thái ta được đường dây, và thành phần hướng được xác định duy nhất bởi hai ô liên tiếp. Chi phí
  mỗi cung \eqref{eq:t5-state} bằng đúng số hạng tương ứng của \eqref{eq:t5-cost}.
\end{loigiai}

Mọi trạng thái của ô ra đều là đích (tập \ma{isGoal}), vì đi tới ô ra từ hướng nào cũng được.
Manhattan nhất quán vì mỗi bước tốn ít nhất $1$ và làm khoảng cách Manhattan giảm tối đa $1$;
Chebyshev $\le$ Manhattan nên cũng nhất quán nhưng kém thông tin hơn. Đồ thị trạng thái có
$5mn$ đỉnh và tối đa $20mn$ cung, nên độ phức tạp vẫn là $O(mn \log(mn))$.

\section{Vết thực thi trên bảng mạch mẫu}
\label{sec:t5-vet}

Đi dây từ $(0,0)$ tới $(3,3)$ trên \Hinh{fig:t5-luoi}, \ma{mode} $= 1$. Trạng thái được viết
$(x,y)^{d}$ với $d \in \{U, D, L, R, -\}$.

\begin{bangbieu}{Vết A* của Task 5 trên đồ thị trạng thái (Manhattan)}{tab:t5-vet}{C{0.8cm} C{3.2cm} Y}
  \hangtieude \thead{Vòng} & \thead{Mở rộng $(g; h; f)$} & \thead{Đưa vào \OPEN} \\
  \midrule
  1 & $(0,0)^{-}\ (0; 6; 6)$ & $(1,0)^{D}\ (1; 5; 6)$, $(0,1)^{R}\ (1; 5; 6)$ \\
  2 & $(1,0)^{D}\ (1; 5; 6)$ & $(0,0)^{U}\ (2{,}5; 6; 8{,}5)$ (quay đầu bị phạt), $(2,0)^{D}\ (2; 4; 6)$ \\
  3 & $(2,0)^{D}\ (2; 4; 6)$ & $(1,0)^{U}\ (3{,}5; 5; 8{,}5)$, $(2,1)^{R}\ (3{,}5; 3; 6{,}5)$ \textit{(bẻ góc)} \\
  4 & $(0,1)^{R}\ (1; 5; 6)$ & $(0,0)^{L}\ (2{,}5; 6; 8{,}5)$, $(0,2)^{R}\ (2; 4; 6)$ \\
  5 & $(0,2)^{R}\ (2; 4; 6)$ & $(0,1)^{L}\ (3{,}5; 5; 8{,}5)$, $(0,3)^{R}\ (3; 3; 6)$ \\
  6 & $(0,3)^{R}\ (3; 3; 6)$ & $(1,3)^{D}\ (4{,}5; 2; 6{,}5)$ \textit{(bẻ góc)}, $(0,2)^{L}\ (4{,}5; 4; 8{,}5)$ \\
  7 & $(1,3)^{D}\ (4{,}5; 2; 6{,}5)$ & $(0,3)^{U}\ (6; 3; 9)$, $(2,3)^{D}\ (5{,}5; 1; 6{,}5)$ \\
  8 & $(2,3)^{D}\ (5{,}5; 1; 6{,}5)$ & $(1,3)^{U}\ (7; 2; 9)$, $(3,3)^{D}\ (6{,}5; 0; 6{,}5)$, $(2,2)^{L}\ (7; 2; 9)$ \\
  9 & $(3,3)^{D}\ (6{,}5; 0; 6{,}5)$ & Trạng thái của ô ra được lấy ra, dừng \\
\end{bangbieu}

Ở vòng 4, $(2,1)^{R}$ với $f = 6{,}5$ và $(0,1)^{R}$ với $f = 6$ cùng nằm trong \OPEN: hai nhánh
xuống--phải và phải--xuống có cùng độ dài, nhưng nhánh đi xuống phải bẻ góc sớm vì ô $(3,0)$ bị
chiếm và sẽ phải bẻ góc thêm lần nữa ở $(2,3)$. Nhánh phải--xuống chỉ bẻ góc một lần tại $(0,3)$.

\begin{ketqua}[Kết quả Task 5 trên bảng mạch mẫu]
  Đường dây $(0,0) \to (0,1) \to (0,2) \to (0,3) \to (1,3) \to (2,3) \to (3,3)$: dài 6 đơn vị,
  một lần bẻ góc, $g = 6 + 0{,}5 = 6{,}5$. Đường qua $(2,0), (2,1), (2,2), (2,3)$ cũng dài 6 nhưng
  bẻ góc hai lần ($g = 7$) nên bị loại.
\end{ketqua}

\section{So sánh heuristic}
\label{sec:t5-so-sanh}

\begin{bangbieu}{Manhattan và Chebyshev trên bảng mạch ngẫu nhiên, 20\% ô bị chiếm (chi phí / số trạng thái mở rộng)}{tab:t5-so-sanh}{C{2.4cm} Z Z}
  \hangtieude \thead{Kích thước} & \thead{Manhattan} & \thead{Chebyshev} \\
  \midrule
  $8 \times 10$  & \otot 18,5 / 42 & 18,5 / 102 \\
  $9 \times 10$  & \otot 18,0 / 32 & 18,0 / 115 \\
  $10 \times 10$ & \otot 20,0 / 31 & 20,0 / 90 \\
\end{bangbieu}
\nguonbang{\ma{make report-data}, hạt giống 5. Cả hai chế độ đều trùng Dijkstra trên đồ thị trạng thái (kiểm thử 30 bảng mạch).}

Trên lưới 4 hướng, Manhattan là cận dưới chặt của độ dài dây nên mở rộng ít hơn Chebyshev 2,4 tới
3,6 lần mà vẫn tối ưu; đây là lựa chọn phù hợp cho bài toán đi dây.
